\small
\setlength{\tabcolsep}{3pt}
\renewcommand{\arraystretch}{1.15}
\begin{longtable}{p{0.18\textwidth} p{0.22\textwidth} p{0.50\textwidth}}
\caption{Categorias de métodos de filtragem temporal e artigos associados.} \label{tab:filter-categories} \\
\toprule
Definições & Métodos a serem contabilizados & Título dos artigos dos métodos a serem contabilizados \\
\midrule
\endfirsthead
\caption[]{Categorias de métodos de filtragem temporal e artigos associados.} \\
\toprule
Definições & Métodos a serem contabilizados & Título dos artigos dos métodos a serem contabilizados \\
\midrule
\endhead
\midrule
\multicolumn{3}{r}{Continued on next page} \\
\midrule
\endfoot
\bottomrule
\endlastfoot
filtros digitais clássicos & bandpass (anti-aliasing?); bandpass (butterworth); bandpass (butterworth, busca exaustiva); bandpass (chebyshev i); bandpass (chebyshev ii); bandpass (chebyshev); bandpass (elliptic); bandpass (equiripple); bandpass (fir); bandpass (fir anti-aliasing); bandpass (fir, coefficient decimation); bandpass (fir, constrained least squares); bandpass (fir, hamming window); bandpass (iir); bandpass (indeterminado); bandpass (kaiserwin); bandpass (notch); bandpass (savitzky-golay); bandstop (comb notch); bandstop (indeterminado); bandstop (notch butterworth 5th ordem); bandstop (notch); filtro proprio (morphology); high-pass (butterworth 4th ordem); high-pass (fir); high-pass (indeterminado); high-pass (intederminado); low-pass (butterworth 4th ordem); low-pass (butterworth); low-pass (indeterminado); morphological multiresolution analysis (haar-type morphology) & Common spatio-time-frequency patterns for motor imagery-based brain machine interfaces\newline A hybrid EEG-fNIRS BCI: Motor imagery for EEG and mental arithmetic for fNIRS\newline Motor imagery classification by means of source analysis for brain-computer interface applications\newline A Graph-Based Feature Extraction Algorithm Towards a Robust Data Fusion Framework for Brain-Computer Interfaces\newline Design method of morphological structuring function for pattern recognition of EEG signals during motor imagery and cognition\newline A2FWPO: Anti-aliasing Filter Based on Whale Parameter Optimization Method for Feature Extraction and Recognition of Dance Motor Imagery EEG\newline Multivariate Adaptive Autoregressive Modeling and Kalman Filtering for Motor Imagery BCI\newline Hybrid mental tasks based human computer interface via integration of pronunciation and motor imagery\newline Wavelet analysis of EEG signals during motor imagery\newline Time-frequency analysis of EEG signals from healthy subjects allocated by gender for a subject-independent BCI-based on motor imagery\newline Feature selection strategy for classification of single-trial EEG elicited by motor imagery\newline A Deep Learning Scheme for Motor Imagery Classification based on Restricted Boltzmann Machines\newline Statistical Wavelets With Harmony Search- Based Optimal Feature Selection of EEG Signals for Motor Imagery Classification\newline Emotion-Inducing Imagery Versus Motor Imagery for a Brain-Computer Interface\newline Classification scheme for arm motor imagery\newline Adaptive modelling of eeg signals to produce accurate time-frequency decompositions for use in BCI\newline A study on the efficient tracking of ERD based on the adaptive identification of the subject's reactive band\newline Effects of Trial-Adjusted Neurofeedback Training on Motor-Imagery Based Brain-Computer Interface Performance\newline An Effect of Limb Position in Motor Imagery Training Paradigm in Immersive Virtual Environment\newline Motor-Related EEG Analysis Using a Pole Tracking Approach\newline Enhanced BCI Performance using Diffusion Model for EEG Generation\newline Simultaneous design of FIR filter banks and spatial patterns for EEG signal classification\newline Research on Motor Imagery Signal Recognition Method Based on Wavelet Packet Transform and ConvLSTM\newline Trust and explainability in robotic hand control via adversarial multiple machine learning models with EEG sensor data fusion: A fuzzy decision-making solution\newline A Region of Interest-Based Electrophysiological Source Imaging Technology and its Applications in Analysis of Motor Imagery EEG Signals\newline Classification of EEG Signals Using a Common Spatial Pattern Based Motor-Imagery for a Lower-limb Rehabilitation Exoskeleton\newline State-of-the-Art Versus Deep Learning: A Comparative Study of Motor Imagery Decoding Techniques\newline Optimal design of a bank of spatio-temporal filters for EEG signal classification\newline Smart thoughts: BCI based system implementation to detect motor imagery movements\newline Discrimination of motor imagery patterns by electroencephalogram phase synchronization combined with frequency band energy\newline Real-Time Motor Imagery-Based Brain–Computer Interface System by Implementing a Frequency Band Selection\newline Classification of motor imagery by means of cortical current density estimation and Von Neumann entropy\newline Research on Movement Intentions of Human's Left and Right Legs Based on Electro-Encephalogram Signals\newline EEG pattern recognition: An efficient improvement combination of ERD/ERS/Laterality features to create a self-paced BCI system\newline Classification of imagery motor EEG data with wavelet denoising and features selection\newline EEG classification of different imaginary movements within the same limb\newline Structure constrained semi-nonnegative matrix factorization for EEG-based motor imagery classification\newline Implementation of nonlinear classifiers for adaptive autoregressive EEG features classification\newline Envelope detection based on online ICA algorithm and its application to motor imagery classification\newline An adaptive EEG filtering approach to maximize the classification accuracy in motor imagery\newline Real-time single channel EEG motor imagery based Brain Computer Interface\newline Joint spatial-temporal filter design for analysis of motor imagery EEG\newline Motor imagery classification by means of source analysis methods\newline Study of On-Line Adaptive Discriminant Analysis for EEG-Based Brain Computer Interfaces\newline Classifying Single-Trial EEG During Motor Imagery by Iterative Spatio-Spectral Patterns Learning (ISSPL)\newline A Study on the impact of spectral variability in brain-computer interface\newline Novel feature generation and classification for a 2-state Self-paced Brain Computer Interface system\newline Spatio-Temporal Joint Filtering for Transfer Learning in Brain-Computer Interface of Motion Imaginery\newline Improving classification accuracy for motor imagery BCI using DWT based Features\newline Development of an Adaptive Artifact Subspace Reconstruction Based on Hebbian/Anti-Hebbian Learning Networks for Enhancing BCI Performance\newline Features Extraction Technique Analysis of EEG Signal for Motor Imagery Classification\newline BCI-Based Robotic Arm for Medical Rehabilitation\newline Rapid prototyping of an EEG-based brain-computer interface (BCI)\newline Design of a Brain Computer Interface system based on electroencephalogram(EEG)\newline The Influence of Smoothing Filtering Methods on the Performance of an EEG-Based Brain–Computer Interface\newline A Hardware/Software Prototype of EEG-based BCI System for Home Device Control\newline Feature Extraction Techniques for the Classification of Four-Class Motor Imagery Based EEG Data: A Comparison\newline Model and Data Dual-Driven Double-Point Observation Network for Ultra-Short MI EEG Classification\newline Classification of motor imagery BCI using multivariate empirical mode decomposition\newline A Novel Classification Model Based on DWT and CNN-LSTM Motor EEG Imagination Signals\newline A tensor-based scheme for stroke patients' motor imagery EEG analysis in BCI-FES rehabilitation training\newline Graph-based feature learning methods for subject-dependent and subject-independent motor imagery EEG decoding\newline Time window and frequency band optimization using regularized neighbourhood component analysis for Multi-View Motor Imagery EEG classification\newline A novel simplified convolutional neural network classification algorithm of motor imagery EEG signals based on deep learning\newline Classification of single trial motor imagery EEG recordings with subject adapted non-dyadic arbitrary time-frequency tilings\newline Recognition of motor imagery tasks for BCI using CSP and chaotic PSO twin SVM\newline Optimum spatio-spectral filtering network for brain-computer interface\newline A space-time-frequency analysis approach for the classification motor imagery EEG recordings in a Brain Computer Interface Task\newline A boosting-based spatial-spectral model for stroke patients' EEG analysis in rehabilitation training\newline A CNN-based modular classification scheme for motor imagery using a novel EEG sampling protocol suitable for IoT healthcare systems\newline Data-efficient motor imagery decoding in real-time for the cybathlon brain-computer interface race\newline Wavelet Transform Time-Frequency Image and Convolutional Network-Based Motor Imagery EEG Classification\newline Spatiotemporal source tuning filter bank for multiclass EEG based brain computer interfaces\newline EEG pattern recognition based on dual-tree complex wavelet transform and particle swarm optimization\newline Applying Kalman filter in EEG-Based Brain Computer Interface for Motor Imagery classification\newline The offline feature extraction of four-class motor imagery EEG based on ICA and Wavelet-CSP\newline Spatio-temporally regularized common spatial patterns (STR-CSP) for single-trial EEG classification\newline A bacterial foraging optimization and learning automata based feature selection for motor imagery EEG classification\newline Spatial Filter Bank based on Probabilistic Common Spatial Pattern in Motor Imagery BCI systems\newline A Unified Fisher’s Ratio Learning Method for Spatial Filter Optimization\newline Feature extraction by common spatial pattern in frequency domain for motor imagery tasks classification\newline Data adaptive filtering approach to improve the classification accuracy of motor imagery for BCI\newline Wavelet coherence based channel selection for classifying single trial motor imagery\newline A self produced mother wavelet feature extraction method for motor imagery brain-computer interface\newline Classification of Motor Imagery EEG Based on a Time-Frequency Analysis and Second-Order Blind Identification\newline Using Multiresolution Space-Time-Frequency Features for the Classification Motor Imagery Eeg Recordings in a Brain Computer Interface Task\newline Motor imagery based Brain Computer Interface with subject adapted time-frequency tiling\newline Application of neural network to brain-computer interface\newline A novel EMD-based Common Spatial Pattern for motor imagery brain-computer interface\newline Comparative study of band-power extraction techniques for Motor Imagery classification\newline Data-Driven Frequency Bands Selection in EEG-Based Brain-Computer Interface\newline Efficient EEG-Based Motor Imagery Classification via Wavelet Genetic Optimization and Machine Learning\newline Deep convolutional neural network for decoding motor imagery based brain computer interface\newline Common Spatial Pattern and Riemannian Manifold-Based Real-Time Multiclass Motor Imagery EEG Classification\newline Classification of Left/Right Hand Movement from EEG Signal by Empirical Mode Decomposition with Ensemble of classifiers\newline Classification of Motor Imagery EEG Signals Using Deep Learning\newline Time-Frequency-Space EEG Decoding Model Based on Dense Graph Convolutional Network for Stroke\newline Brain-Computer Interface Based Control for Disabled\newline Motor Imagery EEG Recognition Based on Weight-Sharing CNN-LSTM Network\newline Employing Deep Learning and Discrete Wavelet Transform Approach to Classify Motor Imagery Based Brain Computer Interface System\newline An Algorithm for Motor Imagery Classification Based on Transfer Learning and Feature Fusion\newline Motor Imagery EEG Classification Based on Riemannian Sparse Optimization and Dempster-Shafer Fusion of Multi-Time-Frequency Patterns\newline Enhancement of Time-Frequency Energy for the Classification of Motor Imagery Electroencephalogram Based on an Improved FitzHugh–Nagumo Neuron System\newline Motor Imagery BCI Classification Using Synchrosqueezing Transform\newline Sparse Bayesian Learning for End-to-End EEG Decoding\newline EEGNet-MSD: A Sparse Convolutional Neural Network for Efficient EEG-Based Intent Decoding\newline Motor Imagery Classification Enhancement using Generative Adversarial Networks for EEG Spectrum Image Generation\newline Motor Imagery EEG Classification Using Fine-Tuned Deep Convolutional EfficientNetB0 Model\newline Multi-Task EEG Signal Classification Using Correlation-Based IMF Selection and Multi-Class CSP\newline Enhancing Motor Imagery EEG Signal Classification with Simplified GoogLeNet\newline TMAnet: An Innovative Approach for Classifying Motor Imagery EEG Signals Related to Motor Movements\newline A Novel Motor Imagery EEG Decoding Method Integrating Cascaded Common Spatial Pattern and Riemannian Geometry\newline Hybrid CNN-LSTM Model for EEG-Based Motor Imagery Classification\newline Integrating Virtual Reality and Multifeature Fusion Deep Learning to Improve EEG-Based Motion Imagery Decoding\newline Improving Motor Imagery EEG Signal Classification Using Deep Learning Techniques\newline BCI for Mobile Devices: Time–Frequency and Representation Learning Analysis of Mobile Gesture Tasks\newline Decoding of Intuitive Motor Imagery Paradigm Based on Frequency-Band Double-Branch Network\newline Scale Specific Riemannian Features with Graph Convolutional Network for Motor Imagery EEG Classification\newline SSTAF: Spatial-Spectral-Temporal Attention Fusion Transformer for Motor Imagery Classification\newline Home Robot Motor Imagery Interaction and Emotional Awareness Based on Multimodal EEG Feature Extraction\newline A Novel Deep Learning Framework for High-Gamma EEG Signal Classification Using Hybrid Compact-CNN-LSTM Network\newline Baseline-Guided Representation Learning for Noise-Robust EEG Signal Classification\newline A Novel Motor Imagery Decoding Method Based on MODWT and Light-Weight CNN\newline Brain Computer Interfaces for Enhancing Motor Imagery Classification Using EEG Signal\newline A Classification and Recognition Method for Motor Imagery EEG Signals Using Fusion of DRSN and WPD\newline SKSP: Selected Kernel Spatial Pattern Analysis in Euclidean and Hilbert Space for Decoding Motor Imagery\newline Spatiotemporal Filtering Bank for Decoding Motor Imagery EEG Signals\newline Explainable Common Spatial Pattern to Improve the Design of Brain-Computer Interfaces\newline An Adaptive Convolutional Neural Network With Spatio-Temporal Attention and Dynamic Pathways (ACNN-STADP) for Robust EEG-Based Motor Imagery Classification\newline EEG Processing for Brain-Computer Interface using Wavelet-Decomposition Ant Colony Optimization and Machine Learning\newline Decoding Two-Class Motor Imagery EEG with Capsule Networks\newline Spatial Component-wise Convolutional Network (SCCNet) for Motor-Imagery EEG Classification\newline Frequency-Optimized Local Region Common Spatial Pattern Approach for Motor Imagery Classification\newline A Deep Learning Framework for Decoding Motor Imagery Tasks of the Same Hand Using EEG Signals\newline Enhancing Mu-based BCI Performance with Rhythmic Electrical Stimulation at Alpha Frequency\newline A Novel BCI System Based on Hybrid Features for Classifying Motor Imagery Tasks\newline Motor Imagery EEG Signal Processing and Classification Using Machine Learning Approach\newline EEG-based BCI system for classifying motor imagery tasks of the same hand using empirical mode decomposition\newline A Comparative Study of Wavelet and CSP Features Classified Using LDA, SVM and ANN in EEG Based Motor Imagery\newline Study of wavelet-based performance enhancement for motor imagery brain-computer interface\newline Sensorimotor Cortex EEG signal classification using Hidden Markov Models and Wavelet Decomposition\newline Deep Fusion Feature Learning Network for MI-EEG Classification\newline DeepMI: Deep Learning for Multiclass Motor Imagery Classification\newline Bayesian Optimized Spectral Filters Coupled With Ternary ECOC for Single-Trial EEG Classification\newline Motor Imagery Classification Using TSK Fuzzy Inference Neural Networks\newline Recurrent Deep Learning for EEG-based Motor Imagination Recognition\newline Transform based feature construction utilizing magnitude and phase for convolutional neural network in EEG signal classification\newline Sparse CSP Algorithm via Joint Spatio-Temporal Filtering\newline Time-frequency decomposition-based weighted ensemble learning for motor imagery EEG classification\newline Short-Time Fourier Transform Covariance and Selection, A Feature Extraction Method for Binary Motor Imagery Classification\newline Motor Imagery Tasks EEG Signals Classification Using ResNet with Multi-Time-Frequency Representation\newline A Power Spectrum Pattern Difference-Based Time-Frequency Sub-Band Selection Method for MI-EEG Classification\newline Feature Extraction of EEG Signal by Power Spectral Density for Motor Imagery Based BCI\newline The Classification of Motor Imagery EEG Signals Based on the Time-Frequency-Spatial Feature\newline Deep Temporal-Spatial Feature Learning for Motor Imagery-Based Brain–Computer Interfaces\newline EEG Classification for Motor Imagery Mental Tasks Using Wavelet Signal Denoising\newline Spectrum-Weighted Tensor Discriminant Analysis for Motor Imagery-Based BCI\newline Canonical correlation analysis of EEG for classification of motor imagery\newline Learning Common Time-Frequency-Spatial Patterns for Motor Imagery Classification\newline A Brain-computer Interface for Motor Imagination Based on Multi-channel EEG Analysis\newline Trial Regeneration With Subband Signals for Motor Imagery Classification in BCI Paradigm\newline Spatiotemporal-Filtering-Based Channel Selection for Single-Trial EEG Classification\newline Identification of Anisomerous Motor Imagery EEG Signals Based on Complex Algorithms\newline Recognition of motor intentions from EEGs of the same upper limb by signal traceability and Riemannian geometry features\newline Toward optimal feature and time segment selection by divergence method for EEG signals classification\newline Electroencephalogram Feature Extraction and Classification of Motor Imagery EEG Based on Machine Learning Algorithms\newline Motor imagery classification method based on Riemannian Manifold and CNN-LSTM for advanced brain computer interfaces\newline System based on subject-specific bands to recognize pedaling motor imagery: Towards a BCI for lower-limb rehabilitation\newline A novel classification method for EEG-based motor imagery with narrow band spatial filters and deep convolutional neural network\newline Nonnegative matrix factorization common spatial pattern in brain machine interface\newline Three-class classification of motor imagery EEG data including “rest state” using filter-bank multi-class Common Spatial pattern\newline The comparison of automatic artifact removal methods with robust classification strategies in terms of EEG classification accuracy\newline Robot motion control using Brain Computer Interface\newline EEG classification based on variance\newline Classification of four categories of EEG signals based on relevance vector machine\newline BCI system to classify EEG stimulus pattern\newline Shrinkage estimator based regularization for EEG motor imagery classification\newline Classification of EEG signals by ICA and OVR-CSP\newline An Improved CSP Algorithm and Application in Motor Imagery Recognition\newline Multi-frequency band common spatial pattern with sparse optimization in Brain-Computer Interface\newline Design an Accurate BCI System by a Novel Method and the Minimum Number of Channels Based on Five Class Motor Imagery\newline Short-Term Fourier Transform as Preprocessing for Common Spatial Pattern\newline EEG based Motor Imagery of Hands Classification using Modified Sub-Band Common Spatial Pattern\newline Adaptive Spatial Filtering and Linear Discriminant Learning for Robust Motor Imagery EEG Classification\newline Enhanced Brain-Computer Interface for EEG Motor Imagery Using DC-GCN and Evolutionary Optimization\newline Benchmarking motor imagery algorithms for pediatric users of brain-computer interfaces\newline Motor Imagery Classification in EEG signals using Hybrid CNN-ISqueeze net\newline Enhancing Four-Class Motor Imagery Detection Through Advanced Feature Extraction Techniques\newline Common Spatial Pattern Variants for Feature Extraction in Multi-Class Motor Imagery BCI\newline TriNet: A Hybrid Feature Integration Approach for Motor Imagery Classification in Brain-Computer Interface\newline EEG based BCI for Seemless Home Automation\newline BCI based EEG Signals classification using Machine Learning Approach\newline Motor-imagery EEG signal decoding using multichannel-empirical wavelet transform for brain computer interfaces\newline Cognitive Computing for Brain–Computer Interface-Based Computational Social Digital Twins Systems\newline Improved Decoding of EEG-Based Motor Imagery Using Convolutional Neural Network and Data Space Adaptation\newline Subject-Independent Brain–Computer Interfaces Based on Deep Convolutional Neural Networks\newline Motor Imagery based Multivariate EEG Signal Classification for Brain Controlled Interface Applications\newline A New Channel Selection Method using Autoencoder for Motor Imagery based Brain Computer Interface\newline A Long Short-Term Memory Autoencoder Approach for EEG Motor Imagery Classification\newline A Comparative Study of Different Feature Extraction Methods for Motor Imagery EEG Decoding within the Same Upper Extremity\newline A Novel Channel Selection Method for BCI Classification Using Dynamic Channel Relevance\newline CNN-based Approaches For Cross-Subject Classification in Motor Imagery: From the State-of-The-Art to DynamicNet\newline Filter Bank Sinc-ShallowNet with EMD-based Mixed Noise Adding Data Augmentation for Motor Imagery Classification\newline A Brain to UAV Communication Model using Stacked Ensemble CSP algorithm based on Motor Imagery EEG signal\newline Application of Transfer Learning in Optimized Filter-Bank Regularized CSP to Classification of EEG Signals with Small Dataset\newline Brain-Computer Interface Rehabilitation System Design Based on Motor Imagery\newline Motor Imagery EEG Signal Classification for Stroke Survivors Rehabilitation\newline A Tensor-Based Frequency Features Combination Method for Brain–Computer Interfaces\newline Motor Imagery EEG Signal Classification based on Deep Transfer Learning\newline A Computationally Efficient Multiclass Time-Frequency Common Spatial Pattern Analysis on EEG Motor Imagery\newline An Approach of One-vs-Rest Filter Bank Common Spatial Pattern and Spiking Neural Networks for Multiple Motor Imagery Decoding\newline A Novel Multiple Motor Imagery Experimental Paradigm Design and Neural Decoding\newline Attention-based spatio-temporal-spectral feature learning for subject-specific EEG classification\newline Advanced TSGL-EEGNet for Motor Imagery EEG-Based Brain-Computer Interfaces \\
métodos decompositores de sinal & decomposition (local cosine packets); eemd; emd; filterbank (band tree search, 4 bandas); filterbank (butterworth); filterbank (chebshev ii); filterbank (chebyshev ii); filterbank (chebyshev ii, seleciona bandpass); filterbank (chebyshev); filterbank (constant-q scheme); filterbank (fft + inverse fft); filterbank (fir); filterbank (fir, coefficient decimation); filterbank (fir, least-square); filterbank (gabor); filterbank (iir); filterbank (indeterminado); filterbank (indeterminado, seleciona bandas); filterbank (indeterminado, seleciona bandpass); multivariate emd; masking emd; vmd (decomposition); wavelet (dt-cwt); wavelet (dwt); wavelet (ewt); wavelet (mewt, espaco-temporal); wavelet (modwt); wavelet (swt); wavelet (wpd); wavelet (wpt); wavelet (wtt) & Common spatio-time-frequency patterns for motor imagery-based brain machine interfaces\newline A multi-modal BCI system based on motor imagery\newline Analysis of EEG signal evoked by passive movement and motor imagery\newline Hybrid mental tasks based human computer interface via integration of pronunciation and motor imagery\newline Wavelet analysis of EEG signals during motor imagery\newline Feature selection strategy for classification of single-trial EEG elicited by motor imagery\newline An enhanced time-frequency-spatial approach for motor imagery classification\newline Statistical Wavelets With Harmony Search- Based Optimal Feature Selection of EEG Signals for Motor Imagery Classification\newline Emotion-Inducing Imagery Versus Motor Imagery for a Brain-Computer Interface\newline EEG signal processing: Feature extraction, selection and classification methods\newline Flexible coding scheme for robotic arm control driven by motor imagery decoding\newline Classification scheme for arm motor imagery\newline An Effect of Limb Position in Motor Imagery Training Paradigm in Immersive Virtual Environment\newline Simultaneous design of FIR filter banks and spatial patterns for EEG signal classification\newline Research on Motor Imagery Signal Recognition Method Based on Wavelet Packet Transform and ConvLSTM\newline Trust and explainability in robotic hand control via adversarial multiple machine learning models with EEG sensor data fusion: A fuzzy decision-making solution\newline A Region of Interest-Based Electrophysiological Source Imaging Technology and its Applications in Analysis of Motor Imagery EEG Signals\newline Feature bands study of MI signals based on FBCSP for CNN\newline Enhanced Online Continuous Brain-Control by Deep Learning-Based EEG Decoding\newline Adaptive subject-based feature extraction in brain-computer interfaces using wavelet packet best basis decomposition\newline Comparative analysis of spectral approaches to feature extraction for EEG-based motor imagery classification\newline Supervised Relevance Analysis for Multiple Stein Kernels for Spatio-Spectral Component Selection in BCI Discrimination Tasks\newline EEG classification of different imaginary movements within the same limb\newline Two Class Motor Imagery Classification Based on ANFIS\newline An empirical mode decomposition based filtering method for classification of motor-imagery EEG signals for enhancing brain-computer interface\newline Motor imagery classification using feature relevance analysis: An Emotiv-based BCI system\newline Maximum dependency and minimum redundancy-based channel selection for motor imagery of walking EEG signal detection\newline Feature extraction and visualization of MI-EEG by LLE algorithm\newline Enhancement of Performance of an EEG-based Brain-Computer Interface by means of a Time-Frequency Approach\newline Algorithm of Imagined Left-Right Hand Movement Classification Based on Wavelet Transform and AR Parameter Model\newline Classifying motor imagery EEG by Empirical Mode Decomposition based on spatial-time-frequency joint analysis approach\newline Dynamic initiation and dual-tree complex wavelet feature-based classification of motor imagery of swallow EEG signals\newline A motor imagery BCI experiment using wavelet analysis and spatial patterns feature extraction\newline Optimal segments selection for EEG classification\newline Improving classification accuracy for motor imagery BCI using DWT based Features\newline Features Extraction Technique Analysis of EEG Signal for Motor Imagery Classification\newline Image-based Motor Imagery EEG Classification using Convolutional Neural Network\newline Data-driven Data Augmentation for Motor Imagery Brain-Computer Interface\newline Relevance of filter bank common spatial patterns using multiple kernel learning in motor imagery\newline Design of a Brain Computer Interface system based on electroencephalogram(EEG)\newline Feature Extraction Techniques for the Classification of Four-Class Motor Imagery Based EEG Data: A Comparison\newline Classification of motor imagery BCI using multivariate empirical mode decomposition\newline A Novel Classification Model Based on DWT and CNN-LSTM Motor EEG Imagination Signals\newline Graph-based feature learning methods for subject-dependent and subject-independent motor imagery EEG decoding\newline Wavelet based feature extraction for classification of motor imagery signals\newline An improved feature extraction method for self-paced brain-computer interface application\newline Time window and frequency band optimization using regularized neighbourhood component analysis for Multi-View Motor Imagery EEG classification\newline Feature Extraction and Classification of EEG Signal Based on Deep Learning\newline Classification of single trial motor imagery EEG recordings with subject adapted non-dyadic arbitrary time-frequency tilings\newline EEG feature extraction and pattern classification based on motor imagery in brain-computer interface\newline EEG feature extraction and pattern recognition during right and left hands motor imagery in brain-computer interface\newline An adjustment strategy on multi-session EEG data for online left/right hand imagery classification\newline Classification of EEG with structural feature dictionaries in a brain computer interface\newline Optimum spatio-spectral filtering network for brain-computer interface\newline Separable common spatio-spectral patterns for motor imagery BCI systems\newline A space-time-frequency analysis approach for the classification motor imagery EEG recordings in a Brain Computer Interface Task\newline A boosting-based spatial-spectral model for stroke patients' EEG analysis in rehabilitation training\newline Class discrepancy-guided sub-band filter-based common spatial pattern for motor imagery classification\newline EEG pattern recognition based on dual-tree complex wavelet transform and particle swarm optimization\newline Feature extraction method of motor imagery EEG based on DTCWT sample entropy\newline Discrimination of EEG related to motor imagery by combined wavelet energy feature with phase synchronization feature\newline Kernel-based relevant feature extraction to support Motor Imagery classification\newline Motor imagery EEG signal classification on DWT and crosscorrelated signal features\newline Optimising frequency band selection with forward-addition and backward-elimination algorithms in EEG-based brain-computer interfaces\newline Motor imagery EEG signal classification scheme based on entropy of intrinsic mode function\newline Dual tree complex wavelet transform domain based motor imagery movements classification scheme from EEG signals\newline Characterization of EEG signals using wavelet transform for motor imagination tasks in BCI systems\newline The offline feature extraction of four-class motor imagery EEG based on ICA and Wavelet-CSP\newline Spatio-temporally regularized common spatial patterns (STR-CSP) for single-trial EEG classification\newline Motor imagery EEG signal classification scheme based on wavelet domain statistical features\newline A bacterial foraging optimization and learning automata based feature selection for motor imagery EEG classification\newline Classification of motor imagery hands movement using levenberg-marquardt algorithm based on statistical features of EEG signal\newline Multiclass motor imagery classification for BCI application\newline A Unified Fisher’s Ratio Learning Method for Spatial Filter Optimization\newline Classification of motor imagery EEG signals based on wavelet transform and sample entropy\newline A feature extraction technique of EEG based on EMD-BP for motor imagery classification in BCI\newline Automatic feature selection based motor imagery movements detection scheme from EEG signals in the Dual Tree Complex Wavelet Transform domain\newline Combined long short-term memory based network employing wavelet coefficients for MI-EEG recognition\newline Classification of motor imagery EEG patterns and their topographic representation\newline Masking empirical mode decomposition-based hybrid features for recognition of motor imagery in EEG\newline Using Multiresolution Space-Time-Frequency Features for the Classification Motor Imagery Eeg Recordings in a Brain Computer Interface Task\newline Motor imagery based Brain Computer Interface with subject adapted time-frequency tiling\newline Feature Extraction and Classification of Mental EEG for Motor Imagery\newline Novel use of Empirical Mode Decomposition in single-trial classification of motor imagery for use in brain-computer interfaces\newline EEG-based motor imagery classification using wavelet coefficients and ensemble classifiers\newline Time-frequency selection in two bipolar channels for improving the classification of motor imagery EEG\newline Non-homogeneous spatial filter optimization for EEG-based brain-computer interfaces\newline Application of neural network to brain-computer interface\newline A novel EMD-based Common Spatial Pattern for motor imagery brain-computer interface\newline Nonlinear and nonstationary framework for feature extraction and classification of motor imagery\newline Efficient EEG-Based Motor Imagery Classification via Wavelet Genetic Optimization and Machine Learning\newline Common Spatial Pattern and Riemannian Manifold-Based Real-Time Multiclass Motor Imagery EEG Classification\newline MSFNet: A Multi-Scale Space-Time Frequency Fusion Network for Motor Imagery EEG Classification\newline Classification of Left/Right Hand Movement from EEG Signal by Empirical Mode Decomposition with Ensemble of classifiers\newline Deep Neural Network-Based Empirical Mode Decomposition for Motor Imagery EEG Classification\newline STOW-Net: Spatio-Temporal Operation Based Deep Learning Network for Classifying Wavelet Transformed Motor Imagery EEG Signals\newline A Multi-Domain Convolutional Neural Network for EEG-Based Motor Imagery Decoding\newline Employing Deep Learning and Discrete Wavelet Transform Approach to Classify Motor Imagery Based Brain Computer Interface System\newline Multi-class Motor Imagery Classification Method Using One-dimensional Convolutional Neural Networks\newline An Algorithm for Motor Imagery Classification Based on Transfer Learning and Feature Fusion\newline FBMSNet: A Filter-Bank Multi-Scale Convolutional Neural Network for EEG-Based Motor Imagery Decoding\newline Motor Imagery EEG Classification Based on Riemannian Sparse Optimization and Dempster-Shafer Fusion of Multi-Time-Frequency Patterns\newline Efficient Ocular Artifacts Removal From EEG Recordings Using State-Space Model\newline Sparse Bayesian Learning for End-to-End EEG Decoding\newline EEGNet-MSD: A Sparse Convolutional Neural Network for Efficient EEG-Based Intent Decoding\newline Multi-Task EEG Signal Classification Using Correlation-Based IMF Selection and Multi-Class CSP\newline A Novel Motor Imagery EEG Decoding Method Integrating Cascaded Common Spatial Pattern and Riemannian Geometry\newline A Multi-Band Self-Attention Network for Motor Imagery Classification\newline Decoding of Intuitive Motor Imagery Paradigm Based on Frequency-Band Double-Branch Network\newline Home Robot Motor Imagery Interaction and Emotional Awareness Based on Multimodal EEG Feature Extraction\newline A Novel Motor Imagery Decoding Method Based on MODWT and Light-Weight CNN\newline A Classification and Recognition Method for Motor Imagery EEG Signals Using Fusion of DRSN and WPD\newline SKSP: Selected Kernel Spatial Pattern Analysis in Euclidean and Hilbert Space for Decoding Motor Imagery\newline Spatiotemporal Filtering Bank for Decoding Motor Imagery EEG Signals\newline AWKNet: A Lightweight Neural Network for Motor Imagery Electroencephalogram Classification Based on Adaptive Wavelet Transform Kolmogorov-Arnold\newline K-Medoids-based ICA and Wavelet-Enhanced Multi-resolution Convolutional Network for Motor Imagery EEG\newline ESSDM: An Enhanced Sparse Swarm Decomposition Method and Its Application in Multi-Class Motor Imagery-Based EEG-BCI System\newline EEG Processing for Brain-Computer Interface using Wavelet-Decomposition Ant Colony Optimization and Machine Learning\newline MDRNet: Multi-Branch with Different Feature Representations Network for Motor Imagery Classification\newline Motor Imagery Hand Movements Recognition Using SVM Classifier and Genetic Algorithm Optimization\newline Frequency-Optimized Local Region Common Spatial Pattern Approach for Motor Imagery Classification\newline Motor Imagery EEG Signals Classification Based on Mode Amplitude and Frequency Components Using Empirical Wavelet Transform\newline Wavelet transform based approach for EEG feature selection of motor imagery data for braincomputer interfaces\newline A Filtering Method for Classification of Motor-Imagery EEG Signals for Brain-Computer Interface\newline EEG Feature Extraction of Motor Imagery Based on WT and STFT\newline A data driven Information theoretic feature extraction in EEG-based Motor Imagery BCI\newline A Novel BCI System Based on Hybrid Features for Classifying Motor Imagery Tasks\newline Motor Imagery EEG Signals Decoding by Multivariate Empirical Wavelet Transform-Based Framework for Robust Brain–Computer Interfaces\newline Classification of Motor Imagery Tasks using Phase Space Reconstruction and Empirical Mode Decomposition\newline Motor Imagery EEG Signal Processing and Classification Using Machine Learning Approach\newline EEG motor imagery signals classification using maximum overlap wavelet transform and support vector machine\newline EEG-based BCI system for classifying motor imagery tasks of the same hand using empirical mode decomposition\newline Study of electroencephalogram feature extraction and classification of three tasks of motor imagery\newline Time frequency analysis of non-stationary motor imagery EEG signals\newline A Comparative Study of Wavelet and CSP Features Classified Using LDA, SVM and ANN in EEG Based Motor Imagery\newline Classification of motor imagery eeg using wavelet envelope analysis and LSTM networks\newline Sensorimotor Cortex EEG signal classification using Hidden Markov Models and Wavelet Decomposition\newline Deep Fusion Feature Learning Network for MI-EEG Classification\newline Motor Imagery Classification Using TSK Fuzzy Inference Neural Networks\newline Quantum-inspired Evolutionary Algorithm for Feature Selection in Motor Imagery EEG Classification\newline Time-frequency decomposition-based weighted ensemble learning for motor imagery EEG classification\newline Motor Imagery Multi-classification based on Canonical Correlation Analysis Feature Fusion\newline Artificial Intelligence in Brain Computer Interface\newline A Power Spectrum Pattern Difference-Based Time-Frequency Sub-Band Selection Method for MI-EEG Classification\newline The Classification of Motor Imagery EEG Signals Based on the Time-Frequency-Spatial Feature\newline Deep Temporal-Spatial Feature Learning for Motor Imagery-Based Brain–Computer Interfaces\newline Chaotic Quantum-inspired Evolutionary Algorithm: enhancing feature selection in BCI\newline EEG Classification for Motor Imagery Mental Tasks Using Wavelet Signal Denoising\newline EEG-based Motor Imagery Classification through Transfer Learning of the CNN\newline Spectrum-Weighted Tensor Discriminant Analysis for Motor Imagery-Based BCI\newline Automated Ocular Artifacts Identification and Removal from EEG Data Using Hybrid Machine Learning Methods\newline Canonical correlation analysis of EEG for classification of motor imagery\newline Motor Imagery Classification Based on Hybrid Feature Extraction and Deep Neural Network\newline Learning Common Time-Frequency-Spatial Patterns for Motor Imagery Classification\newline A Brain-computer Interface for Motor Imagination Based on Multi-channel EEG Analysis\newline Research on Multi-feature Extraction Method in EEG Signal of Motor Imagination\newline Parallel convolutional neural network and empirical mode decomposition for high accuracy in motor imagery EEG signal classification\newline A comparative analysis of masking empirical mode decomposition and a neural network with feed-forward and back propagation along with masking empirical mode decomposition to improve the classification performance for a reliable brain-computer interface\newline Identification of Anisomerous Motor Imagery EEG Signals Based on Complex Algorithms\newline An adaptive filter bank for motor imagery based Brain Computer Interface\newline Toward optimal feature and time segment selection by divergence method for EEG signals classification\newline Learning Temporal Information for Brain-Computer Interface Using Convolutional Neural Networks\newline Functional Connectivity for Motor Imaginary Recognition in Brain-computer Interface\newline Electroencephalogram Feature Extraction and Classification of Motor Imagery EEG Based on Machine Learning Algorithms\newline Reject option to reduce false prediction rates for EEG-motor imagery based BCI\newline Three-class Motor Imagery Classification Based on FBCSP Combined with Voting Mechanism\newline Motor imagery classification method based on Riemannian Manifold and CNN-LSTM for advanced brain computer interfaces\newline A self-adaptive frequency selection common spatial pattern and least squares twin support vector machine for motor imagery electroencephalography recognition\newline Motor Imagery Analysis Based on Filter Bank Common Spatial Pattern\newline Classification of left and right foot kinaesthetic motor imagery using common spatial pattern\newline A New Discriminative Common Spatial Pattern Method for Motor Imagery Brain—Computer Interfaces\newline A novel classification method for EEG-based motor imagery with narrow band spatial filters and deep convolutional neural network\newline Random Forest and Filter Bank Common Spatial Patterns for EEG-Based Motor Imagery Classification\newline Nonnegative matrix factorization common spatial pattern in brain machine interface\newline Three-class classification of motor imagery EEG data including “rest state” using filter-bank multi-class Common Spatial pattern\newline Developing a novel multi-fusion brain-computer interface (BCI) system with particle swarm optimization for motor imagery task\newline The comparison of automatic artifact removal methods with robust classification strategies in terms of EEG classification accuracy\newline Analytic common spatial pattern and adaptive classification for multiclass motor imagery-based BCI\newline Discriminative channel addition and reduction for filter bank common spatial pattern in motor imagery BCI\newline Bi-Directional Imagined Hand Movement Classification Using Low Cost EEG-Based BCI\newline A Novel Mechanism to Fuse Various Sub-Aspect Brain-Computer Interface (BCI) Systems with PSO for Motor Imagery Task\newline Ensemble methods combination for Motor Imagery tasks in Brain Computer Interface\newline EEG-based multi-class motor imagery classification using variable sized filter bank and enhanced One Versus One classifier\newline Restricted exhaustive search for frequency band selection in motor imagery classification\newline Shrinkage estimator based regularization for EEG motor imagery classification\newline Filter Bank Common Spatial Pattern (FBCSP) in Brain-Computer Interface\newline Robust filter bank common spatial pattern (RFBCSP) in motor-imagery-based brain-computer interface\newline Multi-class filter bank common spatial pattern for four-class motor imagery BCI\newline Sub-band Common Spatial Pattern (SBCSP) for Brain-Computer Interface\newline Multi-frequency band common spatial pattern with sparse optimization in Brain-Computer Interface\newline Filter Bank Feature Combination (FBFC) approach for brain-computer interface\newline Optimizing EEG Signal Classification for Motor Imagery BCIs: FilterBank CSP with Riemannian Manifolds and Ensemble Learning Models\newline Design an Accurate BCI System by a Novel Method and the Minimum Number of Channels Based on Five Class Motor Imagery\newline A WPSD-Based Feature Extraction Method of EEG Signal for Motor Imagination\newline Channel-optimized Local Region Riemannian Approach for Motor Imagery Classification\newline Advancements in Temporal Fusion: A New Horizon for EEG-Based Motor Imagery Classification\newline A Hierarchical Intelligent Rehabilitation Robotic System Based on MI-EEG\newline Short-Term Fourier Transform as Preprocessing for Common Spatial Pattern\newline A study of EEG motor imagery classification algorithm based on WPD-AFBCSP\newline AutoEncoder Filter Bank Common Spatial Patterns to Decode Motor Imagery From EEG\newline Non-Uniform Filter Banks-based Minimum Distance to Riemannian Mean Classification in Motor Imagery BCI\newline EEG based Motor Imagery BCI using Four Class Iterative Filtering \& Four Class Filter Bank Common Spatial Pattern\newline Improved Feature Selection Using Equal Grouping Method for EEG Motor Imagery\newline EEG based Motor Imagery of Hands Classification using Modified Sub-Band Common Spatial Pattern\newline Motor Imagery Recognition Using Filter Bank Common Spatial Pattern and Recurrent Neural Networks\newline Benchmarking motor imagery algorithms for pediatric users of brain-computer interfaces\newline Common Spatial Pattern Variants for Feature Extraction in Multi-Class Motor Imagery BCI\newline TriNet: A Hybrid Feature Integration Approach for Motor Imagery Classification in Brain-Computer Interface\newline SMANet: A Model Combining SincNet, Multi-Branch Spatial—Temporal CNN, and Attention Mechanism for Motor Imagery BCI\newline A Study on the Classification Method of Motor Imagery EEG Signals Based on Band Filter Bank Common Space Pattern and Continuous Wavelet Transforms\newline Motor-imagery EEG signal decoding using multichannel-empirical wavelet transform for brain computer interfaces\newline Cognitive Computing for Brain–Computer Interface-Based Computational Social Digital Twins Systems\newline Improved Decoding of EEG-Based Motor Imagery Using Convolutional Neural Network and Data Space Adaptation\newline Subject-Independent Brain–Computer Interfaces Based on Deep Convolutional Neural Networks\newline A New Channel Selection Method using Autoencoder for Motor Imagery based Brain Computer Interface\newline Classification of EEG for Upper Limb Motor Imagery: An Approach for Rehabilitation\newline Improving the efficiency of an EEG-based brain computer interface using Filter Bank Common Spatial Pattern\newline Filter Bank Regularized Common Spatial Pattern Ensemble for Small Sample Motor Imagery Classification\newline A New Motor Imagery EEG Classification Method FB-TRCSP+RF Based on CSP and Random Forest\newline A Bayesian Framework to Optimize Double Band Spectra Spatial Filters for Motor Imagery Classification\newline Relevant Frequency Band Selection using Sequential Forward Feature Selection for Motor Imagery Brain Computer Interfaces\newline A Comparative Study of Different Feature Extraction Methods for Motor Imagery EEG Decoding within the Same Upper Extremity\newline A High Performance Approach for Classification of Motor Imagery EEG\newline An Experimental Analysis of Motor Imagery EEG Signals Using Feature Extraction and Classification Methodologies\newline FBCSP-based Multi-class Motor Imagery Classification using BP and TDP features\newline EEG Channel Selection Based on Correlation Coefficient for Motor Imagery Classification: A Study on Healthy Subjects and ALS Patient\newline Filter Bank Adversarial Domain Adaptation For Motor Imagery Brain Computer Interface\newline CNN-based Approaches For Cross-Subject Classification in Motor Imagery: From the State-of-The-Art to DynamicNet\newline Correlation Alignment in Filter Bank Riemannian Tangent Space for Motor Imagery Classification\newline Multi-class MI-EEG Classification: Using FBCSP and Ensemble Learning Based on Majority Voting\newline A Brain to UAV Communication Model using Stacked Ensemble CSP algorithm based on Motor Imagery EEG signal\newline Application of Transfer Learning in Optimized Filter-Bank Regularized CSP to Classification of EEG Signals with Small Dataset\newline Multi-classification for EEG Motor Imagery Signals using Auto-selected Filter Bank Regularized Common Spatial Pattern\newline Motor Imagery EEG Signal Classification for Stroke Survivors Rehabilitation\newline A Tensor-Based Frequency Features Combination Method for Brain–Computer Interfaces\newline FBCSP and Adaptive Boosting for Multiclass Motor Imagery BCI Data Classification: A Machine Learning Approach\newline A Computationally Efficient Multiclass Time-Frequency Common Spatial Pattern Analysis on EEG Motor Imagery\newline An Approach of One-vs-Rest Filter Bank Common Spatial Pattern and Spiking Neural Networks for Multiple Motor Imagery Decoding\newline A Novel Multiple Motor Imagery Experimental Paradigm Design and Neural Decoding\newline NAO Robot Limb Control Method Based on Motor Imagery EEG\newline Attention-based spatio-temporal-spectral feature learning for subject-specific EEG classification\newline Multi-task recognition motor imagery based on deep learning \\
mapas tempo-frequência & bmflc (tempo frequencia); cqt (tempo frequencia, constant-q transform); cwd (tempo frequencia, choi-willimas distribution); espectrograma (indeterminado, tempo frequencia); essdm (tempo frequencia); gt (tempo frequencia, gabor transform); hht; hhtf; mst (stockwell transform modificada); mwtft (tempo frequencia); st (stockwell transform); stft; tvar (mwbf-rols, wavelet, tempo frequencia); wavelet (cmw); wavelet (cwt); wavelet (sst); wigner-ville (tempo frequencia) & A Graph-Based Feature Extraction Algorithm Towards a Robust Data Fusion Framework for Brain-Computer Interfaces\newline Time-frequency analysis of EEG signals from healthy subjects allocated by gender for a subject-independent BCI-based on motor imagery\newline State-of-the-Art Versus Deep Learning: A Comparative Study of Motor Imagery Decoding Techniques\newline Comparative analysis of spectral approaches to feature extraction for EEG-based motor imagery classification\newline EEG pattern recognition: An efficient improvement combination of ERD/ERS/Laterality features to create a self-paced BCI system\newline Motor imagery classification using feature relevance analysis: An Emotiv-based BCI system\newline Real-time single channel EEG motor imagery based Brain Computer Interface\newline Classifying motor imagery EEG by Empirical Mode Decomposition based on spatial-time-frequency joint analysis approach\newline BCI-Based Robotic Arm for Medical Rehabilitation\newline Enhancing Motor Imagery Classification in Brain–Computer Interfaces Using Deep Learning and Continuous Wavelet Transform\newline Classification of motor imagery BCI using multivariate empirical mode decomposition\newline A tensor-based scheme for stroke patients' motor imagery EEG analysis in BCI-FES rehabilitation training\newline Classification of four-class motor imagery employing single-channel electroencephalography\newline An improved feature extraction method for self-paced brain-computer interface application\newline Feature Extraction and Classification of EEG Signal Based on Deep Learning\newline A novel simplified convolutional neural network classification algorithm of motor imagery EEG signals based on deep learning\newline Motor imagery EEG discrimination using the correlation of wavelet features\newline A CNN-based modular classification scheme for motor imagery using a novel EEG sampling protocol suitable for IoT healthcare systems\newline Wavelet Transform Time-Frequency Image and Convolutional Network-Based Motor Imagery EEG Classification\newline Kernel-based relevant feature extraction to support Motor Imagery classification\newline Motor imagery movements classification using multivariate EMD and short time Fourier transform\newline A self produced mother wavelet feature extraction method for motor imagery brain-computer interface\newline Classification of Motor Imagery EEG Based on a Time-Frequency Analysis and Second-Order Blind Identification\newline BMFLC with neural network and DE for better event classification\newline Application of neural network to brain-computer interface\newline Comparative study of band-power extraction techniques for Motor Imagery classification\newline Novel feature of the EEG based motor imagery BCI system: Degree of imagery\newline Deep convolutional neural network for decoding motor imagery based brain computer interface\newline Optimal Fuzzy Logic Enabled EEG Motor Imagery Classification for Brain Computer Interface\newline A Subject-Specific Time Window Selection Method for Motor Imagery BCI\newline Classification of Motor Imagery EEG Signals Using Deep Learning\newline Classification Algorithm for Motor Imagery EEG Signals Based on Multi-Domain Information CNN\newline Time-Frequency-Space EEG Decoding Model Based on Dense Graph Convolutional Network for Stroke\newline Subject Wise Motor Imagery Classification from EEG Data Using Transfer Learning\newline Motor Imagery based EEG Signal Classification using Multi-scale CNN Architecture\newline Brain-Computer Interface Based Control for Disabled\newline Motor Imagery EEG Recognition Based on Weight-Sharing CNN-LSTM Network\newline Time-frequency analysis and spiking neural network for motor imagery EEG classification\newline Enhancement of Time-Frequency Energy for the Classification of Motor Imagery Electroencephalogram Based on an Improved FitzHugh–Nagumo Neuron System\newline EEG-Channel-Temporal-Spectral-Attention Correlation for Motor Imagery EEG Classification\newline Motor Imagery BCI Classification Using Synchrosqueezing Transform\newline Motor Imagery Classification Enhancement using Generative Adversarial Networks for EEG Spectrum Image Generation\newline Motor Imagery EEG Classification Using Fine-Tuned Deep Convolutional EfficientNetB0 Model\newline Enhancing Motor Imagery EEG Signal Classification with Simplified GoogLeNet\newline Integrating Virtual Reality and Multifeature Fusion Deep Learning to Improve EEG-Based Motion Imagery Decoding\newline BCI for Mobile Devices: Time–Frequency and Representation Learning Analysis of Mobile Gesture Tasks\newline Scale Specific Riemannian Features with Graph Convolutional Network for Motor Imagery EEG Classification\newline SSTAF: Spatial-Spectral-Temporal Attention Fusion Transformer for Motor Imagery Classification\newline Bidirectional Time-Frequency Pyramid Network for Enhanced Robust EEG Classification\newline A Novel Deep Learning Framework for High-Gamma EEG Signal Classification Using Hybrid Compact-CNN-LSTM Network\newline Brain Computer Interfaces for Enhancing Motor Imagery Classification Using EEG Signal\newline ESSDM: An Enhanced Sparse Swarm Decomposition Method and Its Application in Multi-Class Motor Imagery-Based EEG-BCI System\newline Motor Imagery Hand Movements Recognition Using SVM Classifier and Genetic Algorithm Optimization\newline Convolutional Neural Network Based Approach Towards Motor Imagery Tasks EEG Signals Classification\newline Decoding Two-Class Motor Imagery EEG with Capsule Networks\newline A Deep Transfer Convolutional Neural Network Framework for EEG Signal Classification\newline A Deep Learning Framework for Decoding Motor Imagery Tasks of the Same Hand Using EEG Signals\newline Metrological performance of a single-channel Brain-Computer Interface based on Motor Imagery\newline Classification of EEG Motor Imagery Tasks Using Convolution Neural Networks\newline EEG Feature Extraction of Motor Imagery Based on WT and STFT\newline Classification of Motor Imagery EEG Signals with multi-input Convolutional Neural Network by augmenting STFT\newline A convolution neural networks scheme for classification of motor imagery EEG based on wavelet time-frequecy image\newline Study on Classification of Left-Right Hands Motor Imagery EEG Signals Based on CNN\newline Transform based feature construction utilizing magnitude and phase for convolutional neural network in EEG signal classification\newline Short-Time Fourier Transform Covariance and Selection, A Feature Extraction Method for Binary Motor Imagery Classification\newline A Transfer Learning Method based on VGG-16 Convolutional Neural Network for MI Classification\newline Motor Imagery Signal Classification Using Constant-Q Transform for BCI Applications\newline Time-frequency Based EEG Motor Imagery Signal Classification with Deep Learning Networks\newline Motor Imagery Tasks EEG Signals Classification Using ResNet with Multi-Time-Frequency Representation\newline Transfer Learning for Classifying Motor Imagery EEG: A Comparative Study\newline A Power Spectrum Pattern Difference-Based Time-Frequency Sub-Band Selection Method for MI-EEG Classification\newline S-EEGNet: Electroencephalogram Signal Classification Based on a Separable Convolution Neural Network With Bilinear Interpolation\newline Spectrum-Weighted Tensor Discriminant Analysis for Motor Imagery-Based BCI\newline Classification of Motor Imagery EEG Signals Based on Data Augmentation and Convolutional Neural Networks\newline Application of Hilbert-Huang transform for the study of motor imagery tasks\newline System based on subject-specific bands to recognize pedaling motor imagery: Towards a BCI for lower-limb rehabilitation\newline A wavelet-based time-frequency analysis approach for classification of motor imagery for brain-computer interface applications\newline Nonnegative matrix factorization common spatial pattern in brain machine interface\newline A WPSD-Based Feature Extraction Method of EEG Signal for Motor Imagination\newline Short-Term Fourier Transform as Preprocessing for Common Spatial Pattern\newline Motor Imagery Classification in EEG signals using Hybrid CNN-ISqueeze net\newline EEG based BCI for Seemless Home Automation\newline A Study on the Classification Method of Motor Imagery EEG Signals Based on Band Filter Bank Common Space Pattern and Continuous Wavelet Transforms\newline Motor Imagery EEG Signal Classification based on Deep Transfer Learning\newline A Computationally Efficient Multiclass Time-Frequency Common Spatial Pattern Analysis on EEG Motor Imagery\newline An Approach to EEG based BCI for Motor Imagery using Explainable Transfer Learning\newline Performance Evaluation of Xception Networks and Short-Time Fourier Transform Spectrograms for Motor Imagery Classification \\
métodos integrados & integrado; integrado (filterbank, sinc-conv) & A2FWPO: Anti-aliasing Filter Based on Whale Parameter Optimization Method for Feature Extraction and Recognition of Dance Motor Imagery EEG\newline iTa-DFiE: An Innovative Tensor Algebra-Based Detection Framework for Incomplete Noninvasive Electroencephalography\newline CD3Net: A Contrastive Diffusion Model with Domain Adaptive Data Synthetic Network for Motor Imagery and Emotion classification\newline Flexible coding scheme for robotic arm control driven by motor imagery decoding\newline Adaptive spiking convolutional network for decoding motor imagery patterns\newline Enhanced Online Continuous Brain-Control by Deep Learning-Based EEG Decoding\newline EEG denoising with a recurrent quantum neural network for a brain-computer interface\newline Development of an Adaptive Artifact Subspace Reconstruction Based on Hebbian/Anti-Hebbian Learning Networks for Enhancing BCI Performance\newline Model and Data Dual-Driven Double-Point Observation Network for Ultra-Short MI EEG Classification\newline SAT-Net: SincNet-Based Attentive Temporal Convolutional Network for Motor Imagery Classification\newline MSFNet: A Multi-Scale Space-Time Frequency Fusion Network for Motor Imagery EEG Classification\newline A Subject-Specific Time Window Selection Method for Motor Imagery BCI\newline Classification Algorithm for Motor Imagery EEG Signals Based on Multi-Domain Information CNN\newline STOW-Net: Spatio-Temporal Operation Based Deep Learning Network for Classifying Wavelet Transformed Motor Imagery EEG Signals\newline A Multi-Domain Convolutional Neural Network for EEG-Based Motor Imagery Decoding\newline Effect of EOG Artifact Removal on EEG Motor-Imagery Classification\newline FBMSNet: A Filter-Bank Multi-Scale Convolutional Neural Network for EEG-Based Motor Imagery Decoding\newline Enhancement of Time-Frequency Energy for the Classification of Motor Imagery Electroencephalogram Based on an Improved FitzHugh–Nagumo Neuron System\newline Efficient Ocular Artifacts Removal From EEG Recordings Using State-Space Model\newline EEG-Channel-Temporal-Spectral-Attention Correlation for Motor Imagery EEG Classification\newline Sparse Bayesian Learning for End-to-End EEG Decoding\newline EEGNet-MSD: A Sparse Convolutional Neural Network for Efficient EEG-Based Intent Decoding\newline TMAnet: An Innovative Approach for Classifying Motor Imagery EEG Signals Related to Motor Movements\newline Parity FuseNet-Mi: Learning Spatio-Temporal Dynamics via Hierarchical Odd-Even Kernel Convolutions for Motor Imagery Decoding\newline Integrating Virtual Reality and Multifeature Fusion Deep Learning to Improve EEG-Based Motion Imagery Decoding\newline Learned-Weight Ensemble Monte Carlo DropBlock for Uncertainty Estimation and EEG Classification\newline A Multi-Band Self-Attention Network for Motor Imagery Classification\newline Multi-Scale Causal Convolution-Based Swin Transformer for Motor Imagery EEG Classification\newline Decoding of Intuitive Motor Imagery Paradigm Based on Frequency-Band Double-Branch Network\newline SSTAF: Spatial-Spectral-Temporal Attention Fusion Transformer for Motor Imagery Classification\newline RGFTSLANet: A Cross-Subject Classification Model for Decoding EEG-Based Motor Imagery Tasks\newline Home Robot Motor Imagery Interaction and Emotional Awareness Based on Multimodal EEG Feature Extraction\newline Bidirectional Time-Frequency Pyramid Network for Enhanced Robust EEG Classification\newline Baseline-Guided Representation Learning for Noise-Robust EEG Signal Classification\newline A Lightweight End-to-End Three-domain Feature Fusion Network for Motor Imagery Decoding\newline A Novel Motor Imagery Decoding Method Based on MODWT and Light-Weight CNN\newline AWKNet: A Lightweight Neural Network for Motor Imagery Electroencephalogram Classification Based on Adaptive Wavelet Transform Kolmogorov-Arnold\newline K-Medoids-based ICA and Wavelet-Enhanced Multi-resolution Convolutional Network for Motor Imagery EEG\newline An Adaptive Convolutional Neural Network With Spatio-Temporal Attention and Dynamic Pathways (ACNN-STADP) for Robust EEG-Based Motor Imagery Classification\newline Spatial Component-wise Convolutional Network (SCCNet) for Motor-Imagery EEG Classification\newline Transfer Learning based Motor Imagery Classification using Convolutional Neural Networks\newline SincNet-Based Hybrid Neural Network for Motor Imagery EEG Decoding\newline Deep Temporal-Spatial Feature Learning for Motor Imagery-Based Brain–Computer Interfaces\newline S-EEGNet: Electroencephalogram Signal Classification Based on a Separable Convolution Neural Network With Bilinear Interpolation\newline Parallel convolutional neural network and empirical mode decomposition for high accuracy in motor imagery EEG signal classification\newline Classification of Motor Imagery EEG Signals Based on Data Augmentation and Convolutional Neural Networks\newline Learning Temporal Information for Brain-Computer Interface Using Convolutional Neural Networks\newline Motor imagery classification method based on Riemannian Manifold and CNN-LSTM for advanced brain computer interfaces\newline A novel classification method for EEG-based motor imagery with narrow band spatial filters and deep convolutional neural network\newline A Hierarchical Intelligent Rehabilitation Robotic System Based on MI-EEG\newline EEG based Motor Imagery of Hands Classification using Modified Sub-Band Common Spatial Pattern\newline Benchmarking motor imagery algorithms for pediatric users of brain-computer interfaces\newline TriNet: A Hybrid Feature Integration Approach for Motor Imagery Classification in Brain-Computer Interface\newline SMANet: A Model Combining SincNet, Multi-Branch Spatial—Temporal CNN, and Attention Mechanism for Motor Imagery BCI\newline Improved Decoding of EEG-Based Motor Imagery Using Convolutional Neural Network and Data Space Adaptation\newline Filter Bank Adversarial Domain Adaptation For Motor Imagery Brain Computer Interface\newline CNN-based Approaches For Cross-Subject Classification in Motor Imagery: From the State-of-The-Art to DynamicNet\newline Filter Bank Sinc-ShallowNet with EMD-based Mixed Noise Adding Data Augmentation for Motor Imagery Classification\newline Advanced TSGL-EEGNet for Motor Imagery EEG-Based Brain-Computer Interfaces \\
métodos híbridos & hibrido; multivariate emd & Common spatio-time-frequency patterns for motor imagery-based brain machine interfaces\newline Simultaneous design of FIR filter banks and spatial patterns for EEG signal classification\newline Optimal design of a bank of spatio-temporal filters for EEG signal classification\newline Supervised Relevance Analysis for Multiple Stein Kernels for Spatio-Spectral Component Selection in BCI Discrimination Tasks\newline Joint spatial-temporal filter design for analysis of motor imagery EEG\newline Novel feature generation and classification for a 2-state Self-paced Brain Computer Interface system\newline Spatio-Temporal Joint Filtering for Transfer Learning in Brain-Computer Interface of Motion Imaginery\newline Classification of motor imagery BCI using multivariate empirical mode decomposition\newline Spatio-temporally regularized common spatial patterns (STR-CSP) for single-trial EEG classification\newline Data adaptive filtering approach to improve the classification accuracy of motor imagery for BCI\newline Motor imagery movements classification using multivariate EMD and short time Fourier transform\newline A multivariate empirical mode decomposition based filtering for subject independent BCI\newline Deep Neural Network-Based Empirical Mode Decomposition for Motor Imagery EEG Classification\newline An Algorithm for Motor Imagery Classification Based on Transfer Learning and Feature Fusion\newline Sparse Bayesian Learning for End-to-End EEG Decoding\newline TMAnet: An Innovative Approach for Classifying Motor Imagery EEG Signals Related to Motor Movements\newline Spatiotemporal Filtering Bank for Decoding Motor Imagery EEG Signals\newline Efficient CSP Algorithm With Spatio-Temporal Filtering for Motor Imagery Classification\newline A Filtering Method for Classification of Motor-Imagery EEG Signals for Brain-Computer Interface\newline A data driven Information theoretic feature extraction in EEG-based Motor Imagery BCI\newline Subject dependent feature extraction method for motor imagery based BCI using multivariate empirical mode decomposition\newline Sparse CSP Algorithm via Joint Spatio-Temporal Filtering\newline Spatiotemporal-Filtering-Based Channel Selection for Single-Trial EEG Classification\newline A Novel Channel Selection Method for BCI Classification Using Dynamic Channel Relevance \\
métodos de interferência & hilbert (ht); smoothing (gaussian, aplicado mapa tempo-frequencia); smoothing (iir); smoothing (iir, media movel); smoothing (media movel exponencial); smoothing (savitzky-golay); wavelet (soft-threshold denoise); wavelet (threshold denoising) & Feature bands study of MI signals based on FBCSP for CNN\newline Classification of imagery motor EEG data with wavelet denoising and features selection\newline Envelope detection based on online ICA algorithm and its application to motor imagery classification\newline The Influence of Smoothing Filtering Methods on the Performance of an EEG-Based Brain–Computer Interface\newline Motor imagery EEG discrimination using the correlation of wavelet features\newline Classification of motor imagery EEG patterns and their topographic representation\newline Application of neural network to brain-computer interface\newline A Novel Motor Imagery EEG Decoding Method Integrating Cascaded Common Spatial Pattern and Riemannian Geometry\newline Motor Imagery EEG Signals Classification Based on Mode Amplitude and Frequency Components Using Empirical Wavelet Transform\newline A data driven Information theoretic feature extraction in EEG-based Motor Imagery BCI\newline Classification of motor imagery eeg using wavelet envelope analysis and LSTM networks\newline Deep Temporal-Spatial Feature Learning for Motor Imagery-Based Brain–Computer Interfaces\newline A comparative analysis of masking empirical mode decomposition and a neural network with feed-forward and back propagation along with masking empirical mode decomposition to improve the classification performance for a reliable brain-computer interface\newline Learning Temporal Information for Brain-Computer Interface Using Convolutional Neural Networks\newline Functional Connectivity for Motor Imaginary Recognition in Brain-computer Interface\newline Analytic common spatial pattern and adaptive classification for multiclass motor imagery-based BCI\newline EEG based Motor Imagery of Hands Classification using Modified Sub-Band Common Spatial Pattern\newline A High Performance Approach for Classification of Motor Imagery EEG\newline FBCSP-based Multi-class Motor Imagery Classification using BP and TDP features\newline Filter Bank Adversarial Domain Adaptation For Motor Imagery Brain Computer Interface\newline Motor Imagery EEG Signal Classification based on Deep Transfer Learning \\
\end{longtable}
